% ============================================================
\section{Experiment 2: Does Encoder Adaptation Help or Harm?}
\label{sec:exp2}
% ============================================================

Experiment 1 leaves five cells with a measured pre-trained advantage.  This
section applies the intervention to them---and, for the drift comparisons of
\S\ref{sec:forecasting:nonmoirai}, to all 31 cells that carry a paired
frozen-encoder run---and asks the question the design actually manipulates:
does allowing the encoder to adapt help or hurt?


% ============================================================
\subsection{What Fine-Tuning Does to the Five Survivors}
\label{sec:forecasting:results}
% ============================================================

The five gate-passing cells are the entire evidence base for a claim about
damage to valuable features, and they do not support one
(Appendix Table~\ref{tab:gate_survivors}).

\textbf{On three of the five, full fine-tuning is a large improvement.} On
Moirai-Base at $h{=}96$ and $h{=}192$ and Moirai-Large at $h{=}96$,
forgetting is $-31.4\%$, $-36.3\%$ and $-42.4\%$ against zero-shot, negative
in 3/3 seeds each, at CKA $0.460$, $0.396$ and $0.626$.  These are the
cells with a demonstrated pre-trained advantage \emph{and} the largest
representational change among the five, and fine-tuning them cuts MSE by
roughly a third.  Freezing the encoder captures only part of that gain
($-9.1\%$, $-16.1\%$, $-21.6\%$), so encoder adaptation is where most of the
improvement comes from.

\textbf{On the remaining two, the mean harm does not survive inspection.}
Moirai-Small/ETTh2 at $h{=}192$ reads forg.$_\text{B}{=}{+}1.99{\pm}3.32$ with 5 of
10 seeds positive---a coin flip on the sign, and the closest the matrix comes to
degradation.  At $h{=}96$ the positive mean ($+9.71{\pm}10.58$, 4 of 10 seeds) is
produced by one seed that diverges under the \emph{frozen} encoder too, so it is an
optimisation failure and not an encoder-drift effect; we report means including it
throughout and decompose it in Appendix~\ref{app:gatecorrection}.

\textbf{Consequence.}  There is no cell in this matrix where we can say that
fine-tuning reliably destroyed a demonstrated pre-trained capability.
\S\ref{sec:analysis:degradation} states that verdict against the three-clause
definition and reports what it costs the diagnostic we proposed.

% ============================================================
\subsection{Sample Size, and What Mitigation Is For}
\label{sec:forecasting:replication}
% ============================================================

Sweeping $n{\in}\{500,\ldots,10\mathrm{k}\}$ on Moirai-Small/ETTh2 at $h{=}96$
---one of the five survivors, so a sweep where a pre-trained advantage does
exist---separates the two readings directly: drift grows monotonically with $n$
(CKA $0.949$ at 500 to $0.518$ at 10k) while forgetting is non-monotonic and
reverses sign, $+14\%$ at 500, $-2.9\%$ by 2k, $-5.3{\pm}2.1\%$ at 10k (7/10
negative). \textbf{If drift magnitude predicted harm these columns would track;
they do not.}  Whatever harm exists is confined to Small at $n{\leq}1$k, while
the largest spread across seeds sits at the $n{=}5$k crossover
(Table~\ref{tab:sample_sweep}, Figure~\ref{fig:dissociation} and
Appendix~\ref{app:n5k_trajectories} for that bimodal cell and the
protocol comparison).  This sweep predates held-out scoring and is scored on one
split throughout, which is why its $n{=}500$ value ($+14.0{\pm}3.9$) differs
from the held-out mean above ($+9.71{\pm}10.58$);
\S\ref{sec:analysis:heldout} is about exactly that gap.

\textbf{What mitigation is for.} Five strategies spanning unconstrained
fine-tuning to full freezing---LoRA, EWC, L2-SP, frozen encoder, strict
freeze---are reported in Appendix~\ref{app:full_tables}, on
Moirai-Small/ETTh2, the one cell where we ran the whole spectrum: weight
anchoring does not help there and LoRA does, but extending LoRA to eight
further value-cells leaves full fine-tuning ahead on four of them
(Appendix~\ref{app:loravaluecells}).  The screen bounds what any of it can be
for: on 27 of the 32 cells we screened there is no demonstrated advantage to
protect, and on three of the remaining five, unconstrained fine-tuning beats
zero-shot by 31--42\% and the frozen encoder by 20--22 points.

% ============================================================
\subsection{No Cell Meets the Degradation Definition}
\label{sec:analysis:degradation}
% ============================================================

Applying the three-clause definition of \S\ref{sec:method} to all 31
intervention cells returns \textbf{zero cells}.  The count comes from
\texttt{scripts/cell\_matrix.py}, not from hand assembly.

Clause~(i) removes 26 of the 31 outright: they have no measured pre-trained
advantage, so a worse post-fine-tuning score is benign replacement rather than
damage.  Of the five that survive it, three fail clause~(ii) in the strongest
possible way---full fine-tuning \emph{improves} MSE by 31--42\% with 3/3 seeds
agreeing.  That leaves two, and both fail the unanimity requirement:
Moirai-Small/ETTh2 at $h{=}192$ has forg.$_\text{B}{=}{+}1.99{\pm}3.32$ with 5
of 10 seeds positive, and at $h{=}96$ the positive mean is a single diverged
seed (\S\ref{sec:forecasting:results}).  The $h{=}192$ cell satisfies
clauses~(i) and~(iii) and is the closest the matrix comes.  Relaxing
unanimity---to a majority, 70\%, 80\%, or a $t$-test, at either gate
threshold---admits at most that one cell (Appendix~\ref{app:defsensitivity}).
Nor does replacing the count with inference: aggregated over the 15 value-cells
with clustered resampling, $\denc{=}{-}4.3$ with CI $[-13.1,{+}4.5]$ and 7 of 15
positive, so the contrast straddles zero rather than favouring either regime
(Appendix~\ref{app:pairedinference}).
Changing the gate's \emph{estimator} instead admits none under seven of the
ladder's eight rungs and one under a seasonal-naive denominator weaker than ours
on that cell, so loosening the baseline is what manufactures degradation cells
(Appendix~\ref{app:baselines}).

\textbf{What we are and are not saying.}  Not that fine-tuning never damages a
pre-trained capability, but that across 32 screened cells on six standard
benchmarks and three backbones we could not construct one instance that
survives its own definition---fixed before the corrected gate was computed.
These benchmarks are the wrong place to look: on no cell does the advantage
survive every admissible rung of an eight-baseline ladder, and on the five that
clear our ridge gate, fine-tuning on 500--1{,}000 samples usually improves on it.
\S\ref{sec:limitations} states what a cell would have to look like to settle
the question.

\textbf{This is a change of conclusion, not of emphasis.}  An earlier version of
this work reported eight degradation cells; all eight fail clause~(i) once the
gate is computed against the regression it is \emph{defined} as.  Every reading
this paper has withdrawn, and why, is in Appendix~\ref{app:corrections}.

